% ============================================================
%  CHAPTER 4 — System Architecture
% ============================================================
\chapter{System Architecture}
\label{ch:architecture}

\section{Overview}

The system follows a \textbf{perceive--reason--act} control loop inspired by the
SayCan \cite{ahn2022saycan} and VILA \cite{lin2024vila} architectures.
At each iteration, the robot observes its environment, reasons about what to do using
a VLM, and executes a physical skill.
This loop repeats until the task is declared complete by the VLM, a maximum number
of steps is reached, or a safety condition is triggered.

Five high-level components collaborate to implement this loop:

\begin{enumerate}
  \item \textbf{EmbodiedAgent} — the top-level orchestrator.
  \item \textbf{VLMReasoner} — the Gemini API client and prompt manager.
  \item \textbf{PerceptionManager} — camera access and YOLO detection.
  \item \textbf{StateManager} — robot state tracking and scene graph.
  \item \textbf{Skill Library} — executable manipulation and motion modules.
\end{enumerate}

\section{Control Loop}

Figure~\ref{fig:control_loop} illustrates one iteration of the control loop.

\begin{figure}[H]
\centering
\begin{tikzpicture}[node distance=1.0cm and 2.2cm]
  % Nodes
  \node[box, minimum width=4cm] (input)   {Voice / Text Input};
  \node[box, minimum width=4cm, below=of input]  (capture) {Capture RGB Image};
  \node[box, minimum width=4cm, below=of capture](detect)  {YOLO-E open-vocab detection};
  \node[box, minimum width=4cm, below=of detect] (state)   {Assemble State};
  \node[box, minimum width=4cm, below=of state]  (vlm)     {Gemini 2.0 Flash\\(VLM Query)};
  \node[box, minimum width=4cm, below=of vlm]    (afford)  {Affordance Check};
  \node[box, minimum width=4cm, below=of afford] (exec)    {Execute Skill};
  \node[box, minimum width=4cm, below=of exec]   (update)  {Update State + TTS};

  % Side boxes
  \node[sbox, right=2.2cm of vlm]   (skills) {\textbf{Skills}\\grab\_object\\handover\\search\_with\_head\\push\_aside\\navigate\_to\\go\_home / wave};
  \node[sbox, left=2.2cm of afford] (safety) {\textbf{Safety}\\Affordance gate\\Failure counter\\Safe mode};

  % Arrows
  \foreach \a/\b in {input/capture, capture/detect, detect/state, state/vlm,
                      vlm/afford, afford/exec, exec/update}
    \draw[arr] (\a) -- (\b);

  \draw[arr] (exec.east) -- (skills.west);
  \draw[arr] (afford.west) -- (safety.east);
  \draw[arr] (safety.north) |- (vlm.west);

  % Loop back
  \draw[arr] (update.east) -- ++(1.5,0) |- (capture.east)
    node[midway, right, font=\small] {next step};

  % Done branch
  \draw[arr] (vlm.south east) -- ++(1.0,-0.5)
    node[right, font=\small] {skill=done → finish};
\end{tikzpicture}
\caption{One iteration of the perceive--reason--act control loop.}
\label{fig:control_loop}
\end{figure}

\section{Split-Process Deployment}

A central architectural constraint is that the Docker container running ROS Melodic
uses \textbf{Python 3.6}, which is incompatible with Ultralytics YOLO11n (requires
Python $\geq$ 3.8) and other modern deep learning libraries.
The solution is a \textbf{split-process architecture} where compute-intensive AI
inference runs on the host machine while all ROS communication runs inside the container.

\begin{figure}[H]
\centering
\begin{tikzpicture}[node distance=0.6cm and 1.5cm]

  % Docker container box
  \begin{scope}
  \node[draw=darkblue, thick, rounded corners=6pt, fill=lightblue!10,
        minimum width=7cm, minimum height=6.5cm, label={[darkblue]above:\textbf{Docker Container (Python 3.6 / ROS Melodic)}}]
        at (0,0) {};
  \end{scope}

  \node[box, minimum width=6cm] (agent)   at (0, 2.3) {EmbodiedAgent};
  \node[box, minimum width=6cm] (percept) at (0, 1.2) {PerceptionManager v2};
  \node[box, minimum width=6cm] (state)   at (0, 0.1) {StateManager};
  \node[box, minimum width=6cm] (vlmbox)  at (0,-1.0) {VLMReasoner};
  \node[box, minimum width=6cm] (vis)     at (0,-2.1) {Web Visualiser (port 8080)};

  % Host box
  \begin{scope}
  \node[draw=palgreen, thick, rounded corners=6pt, fill=palgreen!5,
        minimum width=5cm, minimum height=2.5cm,
        label={[palgreen]above:\textbf{Host Machine (Python 3.10)}}]
        at (6.5, 0.5) {};
  \end{scope}

  \node[sbox, minimum width=4cm] (yolo)   at (6.5, 1.1) {YOLO-E Service (yoloe-26s-seg)\\(port 5001)};
  \node[sbox, minimum width=4cm] (live)   at (6.5, 0.0) {Live Detector\\(port 8081)};

  % Cloud
  \node[draw=gray, dashed, rounded corners=6pt, minimum width=3cm, minimum height=1.0cm,
        align=center, font=\small] (gemini) at (6.5,-2.1) {Gemini 2.0 Flash\\via Eden AI\\(HTTPS)};

  % Arrows
  \draw[arr] (percept.east) -- node[above, font=\tiny]{HTTP POST /detect} (yolo.west);
  \draw[arr, dashed] (vlmbox.east) -- node[below, font=\tiny]{HTTPS REST} (gemini.west);

  % ROS arrows (internal)
  \draw[arr, gray] (agent.south) -- (percept.north);
  \draw[arr, gray] (percept.south) -- (state.north);
  \draw[arr, gray] (agent.west) -- ++(-0.5,0) |- (vis.west);

\end{tikzpicture}
\caption{Split-process deployment architecture.}
\label{fig:deployment}
\end{figure}

\section{ROS Communication Graph}

The agent communicates with the robot hardware exclusively through ROS topics and
action servers. Table~\ref{tab:ros_graph} lists all ROS interfaces used.

\small
\begin{longtable}{p{7cm}lp{3.5cm}}
\caption{ROS topics and action servers used by the system.}
\label{tab:ros_graph} \\
\toprule
\textbf{Name} & \textbf{Dir.} & \textbf{Purpose} \\
\midrule
\endfirsthead
\multicolumn{3}{c}{\small\itshape (continued)} \\
\toprule
\textbf{Name} & \textbf{Dir.} & \textbf{Purpose} \\
\midrule
\endhead
\path{/xtion/rgb/image_raw}                       & Sub  & RGB frames \\
\path{/xtion/depth_registered/image_raw}          & Sub  & Depth frames \\
\path{/xtion/rgb/camera_info}                     & Sub  & Camera intrinsics \\
\path{/joint_states}                              & Sub  & Joint positions \\
\path{/mobile_base_controller/odom}               & Sub  & Base odometry \\
\path{/agent/camera_annotated}                    & Pub  & Annotated image \\
\path{/agent/status}                              & Pub  & JSON status \\
\path{/planning_scene}                            & Pub  & Collision objects \\
\path{/play_motion}                               & Act  & Pre-defined motions \\
\path{/move_group}                                & Act  & MoveIt planning \\
\path{/tts}                                       & Act  & Text-to-speech \\
\path{/head_controller/follow_joint_trajectory}   & Act  & Head pan/tilt \\
\path{/torso_controller/follow_joint_trajectory}  & Act  & Torso lift \\
\path{/move_base}                                 & Act  & Navigation \\
\path{/clear_octomap}                             & Srv  & Reset octomap \\
\bottomrule
\end{longtable}
\normalsize

\section{Data Flow}

Figure~\ref{fig:dataflow} shows the data flow from user input to physical action
for a typical command (\emph{``Grab the bottle''}).

\begin{figure}[H]
\centering
\begin{tikzpicture}[
  dnode/.style={rectangle, rounded corners=3pt, draw=darkblue,
                fill=lightblue!25, minimum width=3.2cm, minimum height=0.6cm,
                font=\small, align=center},
  darr/.style={-{Stealth[length=5pt]}, thick, draw=darkblue},
  dlabel/.style={font=\tiny\itshape, midway, above},
  node distance=0.5cm and 1.8cm
]
  \node[dnode] (cmd)    {``Grab the bottle''};
  \node[dnode, below=of cmd]    (rgb)   {RGB image\\640×480};
  \node[dnode, below=of rgb]    (yolo)  {YOLO-E\\bboxes + scores};
  \node[dnode, below=of yolo]   (depth) {Depth cloud\\back-projected};
  \node[dnode, below=of depth]  (state) {State dict\\gripper/location/objects};
  \node[dnode, below=of state]  (prompt){System + user\\prompt + image};
  \node[dnode, below=of prompt] (json)  {JSON response\\skill + params};
  \node[dnode, below=of json]   (aff)   {Affordance\\check};
  \node[dnode, below=of aff]    (skill) {Skill execution\\(subprocess / action)};
  \node[dnode, below=of skill]  (hw)    {Robot hardware\\(arm, torso, gripper)};

  \foreach \a/\b/\lbl in {
    cmd/rgb/Xtion camera,
    rgb/yolo/HTTP microservice,
    yolo/depth/point cloud,
    depth/state/StateManager,
    state/prompt/VLMReasoner,
    prompt/json/Gemini API,
    json/aff/EmbodiedAgent,
    aff/skill/if OK,
    skill/hw/ROS actions}
  {
    \draw[darr] (\a) -- node[dlabel]{\lbl} (\b);
  }
\end{tikzpicture}
\caption{Data flow from user command to physical robot action.}
\label{fig:dataflow}
\end{figure}

\section{Multi-Step Task Execution}

For tasks requiring more than one skill (\eg ``find the bottle and give it to me''),
the agent runs the perceive--reason--act loop iteratively.
At each step after the first, the prompt to the VLM is augmented with the list of
already-executed skills:

\begin{mdframed}[style=promptbox]
\small
\begin{verbatim}
Original task: "find the bottle and give it to me"
Completed so far: search_with_head, grab_bottle

If the original task is FULLY complete, respond with skill='done'.
Otherwise, what skill should execute next?
\end{verbatim}
\end{mdframed}

The loop terminates when:
\begin{itemize}[noitemsep]
  \item The VLM responds with \texttt{skill="done"} (task complete).
  \item The same skill is selected twice in a row (convergence heuristic).
  \item The maximum step count (default 5) is reached.
  \item A skill fails and safe mode is entered.
\end{itemize}

\section{Configuration and Environment Variables}

\begin{table}[H]
\centering
\caption{Environment variables controlling system behaviour.}
\label{tab:env_vars}
\begin{tabularx}{\linewidth}{lXl}
\toprule
\textbf{Variable} & \textbf{Purpose} & \textbf{Default} \\
\midrule
\texttt{EDENAI\_API\_KEY}        & Eden AI API key (Gemini 2.0 Flash proxy) & — (required) \\
\texttt{GROQ\_API\_KEY}         & Groq API key (LLaMA-3.3-70b LLM)        & — (required) \\
\texttt{YOLO\_SERVICE\_URL}     & URL of YOLO-E HTTP microservice          & \texttt{http://localhost:5001} \\
\texttt{ROBOT\_IP}              & Robot ROS master IP address              & \texttt{10.68.0.1} \\
\texttt{HOST\_IP}               & Workstation IP for ROS communication     & auto-detect \\
\texttt{ROS\_MASTER\_URI}       & ROS master endpoint (set from ROBOT\_IP) & — (required) \\
\texttt{ROS\_IP}                & Local IP for ROS (set from HOST\_IP)     & — (required) \\
\texttt{TARGET\_DESCRIPTION}    & VLM chain-of-thought hint for grasping   & unset \\
\texttt{OBSTACLE\_OBJECTS}      & JSON list of collision objects (MoveIt)  & unset \\
\texttt{GRASP\_DX}              & Grasp offset X override (m)              & from YAML \\
\texttt{GRASP\_DY}              & Grasp offset Y override (m)              & from YAML \\
\texttt{GRASP\_DZ}              & Grasp offset Z override (m)              & from YAML \\
\texttt{PULSE\_SERVER}          & If set, use PulseAudio mic input         & unset \\
\bottomrule
\end{tabularx}
\end{table}

\section{Neuro-Symbolic Rule Engine}
\label{sec:rule_engine}

A key architectural addition beyond a plain perceive-reason-act loop is the
\textbf{SymbolicRuleEngine} — a forward-chaining inference engine embedded
inside the \texttt{StateManager}.
It maintains a typed fact base and fires a set of predefined rules to fixpoint
(at most 20 iterations per step), deriving new symbolic facts from grounded
3D observations.

\subsection{Fact Base}

Facts are stored as Python tuples in a set for O(1) membership testing:
\begin{lstlisting}[language=Python]
('left_of',        'bottle', 'cup')
('reachable',      'bottle')
('can_grasp',      'bottle')
('has_affordance', 'bottle', 'graspable')
('holding',        'gripper', 'bottle')
('bottle_is_open',)            # semantic fact from VLM
\end{lstlisting}

\subsection{Affordance Ontology}

Static per-class affordances determine what actions are physically meaningful:
\begin{center}
\small
\begin{tabular}{ll}
\toprule
\textbf{Object} & \textbf{Affordances} \\
\midrule
bottle  & graspable, drinkable, rollable \\
cup     & graspable, drinkable \\
remote  & graspable, pressable \\
person  & can\_receive, addressable \\
chair   & sittable, pushable \\
table   & surface, placeable\_on \\
\bottomrule
\end{tabular}
\end{center}

\subsection{Forward-Chaining Rules}

\begin{center}
\small
\begin{tabularx}{\linewidth}{lXl}
\toprule
\textbf{Rule} & \textbf{Condition} & \textbf{Conclusion} \\
\midrule
Symmetry      & \texttt{near(A,B)}                                    & \texttt{near(B,A)} \\
Symmetry      & \texttt{left\_of(A,B)}                                & \texttt{right\_of(B,A)} \\
Symmetry      & \texttt{above(A,B)}                                   & \texttt{below(B,A)} \\
Transitivity  & \texttt{left\_of(A,B) $\wedge$ left\_of(B,C)}        & \texttt{left\_of(A,C)} \\
Co-location   & \texttt{on\_surface(A,T) $\wedge$ on\_surface(B,T)}  & \texttt{co\_located(A,B)} \\
Reachability  & 3D distance $<$ 1.20\,m, $|Y|<$ 0.65\,m             & \texttt{reachable(obj)} \\
Graspability  & \texttt{reachable(A) $\wedge$ graspable(A)}          & \texttt{can\_grasp(A)} \\
Handover      & \texttt{holding(gripper,X) $\wedge$ near(robot,P)}   & \texttt{handover\_possible(X,P)} \\
Stacking      & \texttt{above(A,B) $\wedge$ near\_xy(A,B)}           & \texttt{stacked\_on(A,B)} \\
\bottomrule
\end{tabularx}
\end{center}

The derived fact base is serialised and prepended to every VLM system prompt,
giving the model structured symbolic context alongside the raw image:
\begin{center}
\small\ttfamily
can\_grasp(bottle), reachable(bottle), left\_of(bottle,cup), near\_3d(bottle,cup)
\end{center}

\subsection{3D Spatial Relation Thresholds}

All positions are expressed in the \texttt{base\_footprint} frame
(X\,=\,forward, Y\,=\,left, Z\,=\,up), with hysteresis bands to prevent
oscillation between labels:

\begin{center}
\small
\begin{tabular}{llr}
\toprule
\textbf{Relation} & \textbf{Condition} & \textbf{Threshold} \\
\midrule
\texttt{left\_of(A,B)}   & $A_Y > B_Y + \delta$            & $\delta=0.08$\,m \\
\texttt{in\_front\_of}   & $A_X < B_X - \delta$            & $\delta=0.10$\,m \\
\texttt{above(A,B)}      & $A_Z > B_Z + \delta$            & $\delta=0.08$\,m \\
\texttt{near\_3d(A,B)}   & $\|A-B\|_3 < 0.30$\,m          & Euclidean \\
\texttt{on\_table(A)}    & $A_Z < 0.90$\,m                 & height heuristic \\
\texttt{reachable(A)}    & $A_X<1.20$\,m, $|A_Y|<0.65$\,m & arm envelope \\
\bottomrule
\end{tabular}
\end{center}

Self-referential pairs ($A = B$) are explicitly filtered to prevent
nonsensical facts such as \texttt{left\_of(bottle,\,bottle)}.

\subsection{VLM Semantic Enrichment}

In addition to geometric facts, the VLM returns a
\texttt{semantic\_facts} field in its JSON response — visual observations
that geometry cannot compute:
\begin{itemize}[noitemsep]
  \item Object state: \texttt{bottle\_is\_open}, \texttt{cup\_is\_full}
  \item Social cues: \texttt{person\_is\_reaching}, \texttt{person\_is\_looking}
  \item Scene state: \texttt{table\_is\_cluttered}, \texttt{path\_is\_clear}
\end{itemize}
These are merged into the fact base after each VLM call and used in the next
forward-chaining pass, completing the neuro-symbolic loop.
